\documentclass[aspectratio=169,xcolor=table]{beamer}

\usepackage[utf8]{inputenc}
\usepackage[T1]{fontenc}
\usepackage[portuguese]{babel}
\usepackage{lmodern}
\usepackage{booktabs}
\usepackage{graphicx}
\usepackage{tikz}
\usepackage{amsmath}

\usetheme{DCC}

\setbeamertemplate{itemize items}[circle]
\graphicspath{{imgs/}{./imgs/}}

\author[Antoniel Magalhães]{\textbf{Antoniel Magalhães}}
\title{Determinantes de parto cesáreo na Bahia}
\subtitle{Aprendizado supervisionado com SINASC · IA III}
\institute{Universidade Federal da Bahia}
\date{\today}

\begin{document}

\begin{frame}[plain,noframenumbering]
    \titlepage
\end{frame}

\begin{frame}{Agenda}
    \tableofcontents
\end{frame}

\section{Problema}

\begin{frame}{Problema}
    \begin{block}{Pergunta}
        Quais fatores estão associados ao parto cesáreo na Bahia e como isso varia entre municípios?
    \end{block}

    \begin{itemize}
        \item Base: SINASC (DATASUS) — Declarações de Nascidos Vivos
        \item Recorte: Bahia, 2022--2024
        \item Alvo: \texttt{PARTO} (vaginal vs.\ cesárea)
    \end{itemize}
\end{frame}

\section{Pré-processamento}

\begin{frame}{Pré-processamento — fontes}
    \textbf{Três arquivos raw} (\texttt{data/raw/}), um por ano:

    \vspace{0.6em}
    \begin{table}
        \centering
        \begin{tabular}{lrl}
            \toprule
            Arquivo & Registros & Ano \\
            \midrule
            \texttt{sinasc\_ba\_2022.parquet} & 173.824 & 2022 \\
            \texttt{sinasc\_ba\_2023.parquet} & 170.314 & 2023 \\
            \texttt{sinasc\_ba\_2024.parquet} & 160.002 & 2024 \\
            \midrule
            \textbf{Total (3 datasources)} & \textbf{504.140} & \\
            \bottomrule
        \end{tabular}
    \end{table}

    \vspace{0.4em}
    {\small Pipeline: \texttt{ml/preprocess.py} $\rightarrow$ \texttt{sinasc\_ba\_unified.parquet}}
\end{frame}

\begin{frame}{Pré-processamento — rótulo \texttt{PARTO}}
    \textbf{Primeira limpeza:} manter apenas classes válidas do alvo.

    \vspace{0.4em}
    \begin{columns}[T]
        \begin{column}{0.48\textwidth}
            \textbf{Códigos SINASC}
            \begin{itemize}
                \item \texttt{1} — parto vaginal
                \item \texttt{2} — cesárea
                \item \texttt{9} — ignorado
                \item vazio — registro sem rótulo
            \end{itemize}
        \end{column}
        \begin{column}{0.48\textwidth}
            \textbf{Regra (\texttt{filter\_valid\_parto})}
            \begin{itemize}
                \item \textcolor{green!50!black}{Manter} \texttt{1} e \texttt{2}
                \item \textcolor{red!70!black}{Remover} \texttt{9}, vazios e demais inválidos
            \end{itemize}
        \end{column}
    \end{columns}

    \vspace{0.6em}
    \begin{table}
        \centering
        \small
        \begin{tabular}{rrrrrrr}
            \toprule
            Ano & Total & Vaginal & Cesárea & Ignorado & Vazio & Removidos \\
            \midrule
            2022 & 173.824 & 90.536 & 83.141 & 16 & 131 & 147 \\
            2023 & 170.314 & 85.189 & 85.003 & 8 & 114 & 122 \\
            2024 & 160.002 & 75.742 & 84.174 & 5 & 81 & 86 \\
            \midrule
            \textbf{Total} & \textbf{504.140} & 251.467 & 252.318 & 29 & 326 & \textbf{355} \\
            \bottomrule
        \end{tabular}
    \end{table}
\end{frame}

\begin{frame}{Pré-processamento — base unificada}
    \begin{block}{Após unificação e limpeza de rótulos}
        \centering
        \Large 503.785 registros utilizáveis
    \end{block}

    \vspace{0.6em}
    \begin{table}
        \centering
        \begin{tabular}{lr}
            \toprule
            Etapa & Registros \\
            \midrule
            Raw (3 anos) & 504.140 \\
            Removidos (ignorados + vazios) & $-$355 (0{,}07\%) \\
            \midrule
            \textbf{Base unificada} & \textbf{503.785} \\
            \bottomrule
        \end{tabular}
    \end{table}

    \vspace{0.4em}
    \begin{itemize}
        \item Coluna \texttt{y\_cesarea}: \texttt{1} = cesárea, \texttt{0} = vaginal
        \item Balanceamento: $\approx$ 50{,}1\% cesárea · 49{,}9\% vaginal
    \end{itemize}
\end{frame}

\begin{frame}{Obrigado}
    \begin{center}
        \Huge Obrigado!

        \vspace{0.8em}
        \large Perguntas e comentários
    \end{center}
\end{frame}

\end{document}
